\chapter{Patrols}
\label{ch:u-patrol}

\section{Start a patrol}

\begin{description}
  \item[From the H12] \ui{Routes}, pick the route, AUTO (key \key{A}), \ui{Start}.
  \item[From NETRA] \menupath{Patrols › Start a patrol now} with the robot and the route, or a
    schedule in \menupath{Schedules} (days and time). The robot downloads the route and its
    inspection points itself.
\end{description}

The robot walks onto the route, follows it, stops and inspects at every stop point, and ends at
the route's end. During the patrol the H12 and NETRA show its state:

\begin{center}
\begin{tabularx}{\linewidth}{@{}>{\bfseries}p{0.36\linewidth}X@{}}
  \toprule
  State on the H12 & Meaning \\
  \midrule
  Walking onto the route & on its way to the route \\
  Patrolling & following the route \\
  Approaching a stop & slowing down to a stop point or the end \\
  Inspecting & standing at a stop: the cameras at work \\
  Paused & someone paused it, or the e-stop, REMOTE, localization or the battery \\
  Blocked & something is in its way (below) \\
  Off the route & too far from the route to follow it \\
  Route finished & done \\
  \bottomrule
\end{tabularx}
\end{center}

\begin{center}
\includegraphics[width=\linewidth]{follower}
\end{center}

\section{Obstacles}

The robot slows down near an obstacle and stops before it. Then, on its own:

\begin{enumerate}
  \item it steps to the side and passes the obstacle if there is room (path shift);
  \item if not, it tries a short detour around it;
  \item if the way stays blocked, it backs up a little and tries again, and finally reports
    \ui{Blocked} on the H12 and as an alert in NETRA.
\end{enumerate}

Clear the way and it goes on by itself; or take it over in REMOTE.

\section{Pause, take over, end}

\begin{itemize}
  \item \ui{Pause} / \ui{Resume}: the robot stands and waits.
  \item REMOTE (key \key{A}) during a patrol pauses it: drive the robot, then AUTO and
    \ui{Resume}. Bring it back near its route first.
  \item \ui{End}: the patrol stops; the robot stands where it is.
  \item STOP (key \key{B}): the e-stop; after the release the patrol stays paused until
    \ui{Resume}.
\end{itemize}

\section{When the robot does not know where it is}

The top bar shows \ui{No loc} and the patrol pauses (Paused: localization). This happens after
the robot's odometry restarted. Walk the robot back to its start spot, facing the usual way:

\begin{itemize}
  \item with a recorded tag in view, press \ui{Localize}: the robot finds its place from the tag;
  \item without a tag, sit the robot down on the start spot and restart it (or ask the technician
    to restart its service): the start spot becomes its origin again.
\end{itemize}

Then start the patrol again.
